\chapter{Fondamenti di machine learning supervisionato}\label{cap:fondamenti}

% Capitolo 2 — Fondamenti di machine learning supervisionato.

In questo capitolo introduciamo i concetti fondamentali del machine learning
supervisionato. Sono pochi, ma vanno capiti bene perché tutto il resto del
documento ci si appoggia. Come sempre, useremo analogie da programmazione e
collegheremo ogni concetto al codice reale del progetto. I riferimenti
classici per questo capitolo sono \cite{goodfellow2016deep} e
\cite{bishop2006pattern}.

\section{Dati, feature e label}\label{sec:dati-feature-label}

Il punto di partenza di ogni problema di ML sono i \textbf{dati}. Un dataset
è una tabella: se conosci SQL, una riga di dataset è esattamente una riga di
una tabella SQL.

\begin{itemize}
  \item \textbf{Campione} (o riga): una singola osservazione. Nel nostro
        caso, una singola misurazione dei parametri vitali di un paziente.
  \item \textbf{Feature} (o colonne): le variabili che descrivono il
        campione. Nel nostro caso sono sei: pressione sistolica, pressione
        diastolica, frequenza cardiaca, temperatura, saturazione
        dell'ossigeno e glicemia.
  \item \textbf{Label} (o etichetta): la risposta che vogliamo imparare a
        prevedere. Nel nostro caso è la classe di rischio:
        \classe{basso}, \classe{medio} o \classe{alto}.
\end{itemize}

Nel ML \emph{supervisionato} il dataset contiene sia le feature sia le label:
durante il training il modello vede coppie (input, output atteso) e impara a
produrre l'output giusto a partire dall'input. È come studiare con le
soluzioni in fondo al libro: ogni esercizio ha la risposta, e l'obiettivo è
imparare il metodo, non memorizzare le singole risposte.

\begin{esempio}{Una riga del dataset del progetto}
Una riga del dataset sintetico è un vettore di sei numeri nell'ordine
\codice{FEATURE\_ORDER}, più la label assegnata dal teacher:

\begin{lstlisting}[caption={Esempio di campione etichettato (valori illustrativi)}]
# feature: sistolica, diastolica, frequenza, temperatura, SpO2, glicemia
campione = [120.0, 80.0, 75.0, 36.8, 98.0, 95.0]
label    = "basso"
\end{lstlisting}
\end{esempio}

\section{Il modello come funzione parametrizzata}\label{sec:modello-funzione-parametrizzata}

Un modello di ML è, in fondo, una \textbf{funzione parametrizzata}:

\begin{equation}
y = f(x; \theta)
\end{equation}

dove $x$ è l'input (le feature), $y$ è l'output (la label) e $\theta$ è un
vettore di parametri. La forma di $f$ è scelta dal progettista; i valori di
$\theta$ sono sconosciuti e vengono \emph{imparati} dai dati.

L'analogia da programmatore è immediata: pensa a una funzione con parametri
da tarare.

\begin{lstlisting}[caption={Una funzione parametrizzata in Python}]
def classifica(x, theta):
    # la forma e' fissa, i parametri theta sono da tarare
    return combinazione(x, theta)
\end{lstlisting}

Il \textbf{training} è il processo che trova i valori di $\theta$ che fanno
funzionare bene la funzione sui dati. Nel nostro progetto la funzione $f$ è
un MLP a due strati e i parametri $\theta$ sono le matrici di pesi e i
vettori di bias: \codice{W1}, \codice{b1}, \codice{W2}, \codice{b2} (vedi
\file{src/telemedicina\_supervised/ml/mlp.py}). Il training è un loop di
ottimizzazione che aggiorna questi parametri finché la funzione non
approssima bene il teacher (capitolo \ref{cap:training}).

\section{Training vs test}\label{sec:training-test}

Se valutiamo il modello sugli stessi dati con cui l'abbiamo addestrato,
misuriamo solo la sua capacità di \emph{ricordare}, non di \emph{capire}. Per
misurare la generalizzazione serve un insieme di dati che il modello non ha
mai visto durante il training.

Per questo il dataset viene diviso in tre parti:

\begin{itemize}
  \item \textbf{Train}: i dati su cui il modello impara (aggiorna i
        parametri).
  \item \textbf{Validation}: i dati usati per scegliere gli iperparametri e
        fermare il training quando il modello smette di migliorare.
  \item \textbf{Test}: i dati usati una sola volta, alla fine, per misurare
        quanto il modello generalizza su casi mai visti.
\end{itemize}

Nel nostro progetto lo split è congelato: 40\,000 campioni di train, 10\,000
di validation e 20\,000 di test, generati con seed diversi e salvati su
disco (capitolo \ref{cap:dati}). Il test non viene mai toccato durante il
tuning: è il "giudice" che si consulta una sola volta.

\section{Generalizzazione e overfitting}\label{sec:generalizzazione-overfitting}

Il rischio principale del training è l'\textbf{overfitting}: il modello
impara a memoria il training set invece di capire il pattern sottostante.
L'analogia è con lo studente che ripete a memoria le soluzioni degli
esercizi visti: va benissimo sugli esercizi noti, ma fallisce su quelli
nuovi.

Un modello in overfitting ha una loss bassissima sul train ma alta sul test:
ha memorizzato il rumore e i dettagli irrilevanti dei dati di addestramento
\cite{bishop2006pattern}. Al contrario, un modello che generalizza ha
performance simili su train e test.

Come si difende un progetto da questo rischio?

\begin{itemize}
  \item \textbf{Validation set}: si monitora la loss sulla validation durante
        il training; quando smette di scendere, ci si ferma.
  \item \textbf{Early stopping}: nel nostro progetto il training si ferma
        dopo \codice{PAZIENZA} epoche senza miglioramento della validation
        loss (\codice{PAZIENZA = 5} in \file{config.py}).
  \item \textbf{Test separato}: la valutazione finale avviene su dati mai
        visti, una sola volta.
\end{itemize}

\section{Il vocabolario comune}\label{sec:vocabolario-comune}

Prima di proseguire, fissiamo il vocabolario che useremo in tutto il
documento.

\begin{description}
  \item[Dataset:] l'insieme di campioni (righe) con le loro feature.
  \item[Campione:] una singola osservazione, una riga del dataset.
  \item[Feature:] una colonna del dataset, una variabile in ingresso.
  \item[Label:] la risposta attesa per un campione.
  \item[Predizione:] l'output prodotto dal modello per un input.
  \item[Loss:] un numero che misura quanto le predizioni del modello sono
        lontane dalle label attese. Più è bassa, meglio è.
  \item[Epoca:] un passaggio completo del modello su tutto il training set.
  \item[Batch:] un sottoinsieme di campioni usato per un singolo
        aggiornamento dei parametri.
  \item[Iperparametro:] un valore scelto dal progettista prima del training
        (es. numero di neuroni nascosti, learning rate), da non confondere
        con i parametri $\theta$ che il training impara.
\end{description}

\section{Nel progetto}\label{sec:nel-progetto}

Mettiamo insieme i concetti di questo capitolo con i numeri reali del
progetto:

\begin{itemize}
  \item \textbf{Feature}: 6 parametri vitali, nell'ordine ufficiale
        \codice{FEATURE\_ORDER}.
  \item \textbf{Classi}: 3 (\classe{basso}, \classe{medio}, \classe{alto}).
  \item \textbf{Campioni}: 70\,000 totali, divisi in 40\,000 train, 10\,000
        validation e 20\,000 test.
  \item \textbf{Modello}: una funzione parametrizzata $f(x; \theta)$ con
        $\theta$ = pesi e bias di un MLP a due strati.
  \item \textbf{Training}: un loop di ottimizzazione che minimizza la loss
        (cross-entropy) con minibatch SGD e early stopping.
\end{itemize}

\begin{esempio}{I numeri reali del progetto}
La distribuzione delle classi nel dataset congelato (da
\file{data/models/report.json}) è bilanciata: nel train ci sono 14\,703
campioni \classe{basso}, 11\,423 \classe{medio} e 13\,874 \classe{alto}.
Nessuna classe è dominante, e questo aiuta il modello a imparare tutte e tre
le classi.
\end{esempio}

Il prossimo capitolo (\ref{cap:dati}) entra nel dettaglio dei dati: le sei
feature cliniche, perché il dataset è sintetico e come è stato congelato.